problem:  
Let \((M^n,g)\) be a complete, noncompact Riemannian manifold with nonnegative Ricci curvature and positive asymptotic volume ratio
\[
\theta(M,g)=\lim_{r\to\infty}\frac{\operatorname{Vol}(B(p,r))}{\omega_n r^n}>0.
\]
For each basepoint \(p\in M\), consider any pointed measured Gromov–Hausdorff limit of rescalings \((M,r_i^{-2}g,p,\operatorname{Vol}/r_i^n)\to (C(X_p),o_p,\mu_p)\), which is a metric cone with cross-section \(X_p\).  
Assume that for every \(p\), **the tangent cone at infinity is unique**, i.e. all such limits are isometric to the same \(C(X_p)\).

**Problem:**  
> Under this uniqueness assumption, is the map
> \[
> p\longmapsto [X_p]
> \]
> which sends a basepoint to the isometry class of its cross-section \(X_p\), *continuous* with respect to the Gromov–Hausdorff topology on compact spaces?  
> 
> Equivalently, do nearby points in \(M\) necessarily have asymptotic cones that are close (or even isometric) in the Gromov–Hausdorff sense, or can the large-scale geometry of \(M\) “jump” discontinuously with small variations in the basepoint despite uniform curvature and volume-growth control?  

A refined version asks:  
> If not continuous, must the discontinuity set for \(p\mapsto [X_p]\) be of measure zero or contain directions of noncollapsed splitting?  

Why is it a "good" problem:  
This question formalizes a central missing piece in **Ricci limit geometry**—the basepoint regularity of tangent cones at infinity. While Cheeger–Colding theory provides tangent cones up to subsequences and partial regularity, it is completely unknown whether uniqueness implies any *continuous or rigid structure* across basepoints. Proving continuity would elevate the “geometry at infinity” of \((M,g)\) from a pointwise notion to a well-defined, stable invariant of the manifold; disproving it would reveal a fundamental instability in the asymptotic regime of Ricci ≥ 0 geometry.  

The problem bridges **Riemannian curvature theory** (Ricci bounds, volume comparison) and **metric measure limit theory** (Gromov–Hausdorff convergence and RCD spaces), demanding both analytic and geometric insight. It is simply stated but profoundly deep: it asks whether the asymptotic structure of space under fixed curvature control behaves continuously with position—a concept central to geometric analysis and potential future generalizations of geometric stability theorems.